
% !TEX program = pdflatex
\documentclass[11pt,a4paper]{report}

\usepackage[T1]{fontenc}
\usepackage[utf8]{inputenc}
\usepackage{lmodern}
\usepackage{microtype}

\usepackage{amsmath,amssymb,amsthm,bm}
\usepackage{mathtools}
\usepackage{physics}

\usepackage{graphicx}
\usepackage{subcaption}
\usepackage{booktabs}
\usepackage{siunitx}
\usepackage{hyperref}
\usepackage{cleveref}

\usepackage{listings}
\usepackage{xcolor}
\lstset{
  basicstyle=\ttfamily\small,
  keywordstyle=\color{blue!70!black},
  commentstyle=\color{green!40!black},
  stringstyle=\color{red!60!black},
  breaklines=true,
  frame=single,
  columns=fullflexible
}

\title{Notes for NN Mechanical Cloaking Project\\
Rayleigh-wave cloaking: theory, implementation, and loss construction}
\author{(project notes)}
\date{\today}

\begin{document}
\maketitle
\tableofcontents

\chapter{Overview and scope}
These notes summarize the theoretical and numerical ingredients currently used in the NN mechanical cloaking project:
(i) transformation elasticity (with emphasis on the Milton--Briane--Willis framework),
(ii) Cosserat (micropolar) elasticity as a practical gauge choice for Rayleigh-wave cloaking,
(iii) inverse design options (gradient-free and gradient-based),
(iv) weak-form implementation for 2D in-plane elastodynamics (Rayleigh-wave setting) in \texttt{JAX-FEM},
(v) absorbing boundary strategies compatible with real-valued solvers (Rayleigh/Caughey damping layers),
(vi) reference-field computation templates, and
(vii) loss functions for cloaking.

\paragraph{Project code anchor.}
Implementation details referred to in these notes correspond to the current project file \texttt{boundaries.py}
(in particular: \texttt{RayleighCloakProblem}, \texttt{get\_tensor\_map}, \texttt{get\_mass\_map}, \texttt{get\_surface\_maps},
\texttt{cloaking\_loss}, and the helper functions for the transformation and symmetrized stiffness construction).

\chapter{Transformation elasticity theory (Milton et al.)}
\section{Classical linear elastodynamics in the virtual domain}
Let $\Omega_X$ be the \emph{virtual} (reference) domain, with coordinates $\bm X$ and displacement $\bm U(\bm X,t)$.
In the absence of body forces, the Navier equation reads
\begin{equation}
\nabla_X \cdot \bm\sigma(\bm U) = \rho\, \ddot{\bm U}, \qquad
\bm\sigma(\bm U) = \bm C : \bm\varepsilon(\bm U), \qquad
\bm\varepsilon(\bm U) = \tfrac12(\nabla_X \bm U + (\nabla_X \bm U)^{\mathsf T}),
\end{equation}
where $\bm C$ is the (typically symmetric) fourth-order elasticity tensor with minor and major symmetries (for Cauchy elasticity).

\section{Coordinate transformation and gauge}
Let $\bm x = \bm\Xi(\bm X)$ be an invertible mapping from virtual to \emph{physical} coordinates.
Define the deformation gradient $\bm F = \nabla_X \bm x$ and Jacobian $J=\det \bm F$.

A key point in transformation elastodynamics is that the transformed equation depends on the chosen relation
(\emph{gauge}) between virtual and physical displacements. Milton--Briane--Willis (MBW) showed that for general transformations,
preserving the form of the governing equations can introduce effective density anisotropy and additional coupling terms
(Willis-type constitutive structure).

\subsection{Cosserat gauge (form invariance with non-symmetric stress)}
An alternative gauge is the \emph{identity} gauge
\begin{equation}
\bm U(\bm X,t) = \bm u(\bm x,t),
\end{equation}
which (for 2D in-plane transformations used for carpet cloaks) yields a transformed equation of Navier type:
\begin{equation}
\nabla_x \cdot \left(\bm C^{\mathrm{eff}} : \nabla_x \bm u \right) = \rho^{\mathrm{eff}} \ddot{\bm u},
\end{equation}
with
\begin{equation}
C^{\mathrm{eff}}_{ijkl} = J^{-1}\, F_{jJ}\, F_{lL}\, C_{iJkL}, \qquad \rho^{\mathrm{eff}} = \rho\, J^{-1},
\label{eq:cosserat_ceff}
\end{equation}
(using Einstein summation).
Note that $\bm F$ acts on the \emph{derivative} indices $(j,l)$, not the DOF indices $(i,k)$,
because the coordinate transformation enters through the chain rule
$\partial/\partial X_J = F_{jJ}\,\partial/\partial x_j$.

\paragraph{Derivation.}
Starting from the reference weak form
\begin{equation}
\int_{\Omega_0} C_{IJKL}\,\frac{\partial u_K}{\partial X_L}\,\frac{\partial v_I}{\partial X_J}\,d\Omega_0,
\end{equation}
the Cosserat gauge sets $u_i(\bm x) = U_I(\bm X)$ with $i=I$.
Applying the chain rule $\partial/\partial X_J = F_{jJ}\,\partial/\partial x_j$
and the volume element $d\Omega_0 = d\Omega/J$, we obtain
\begin{equation}
\int_{\Omega} \underbrace{J^{-1}\,F_{jJ}\,F_{lL}\,C_{iJkL}}_{= C^{\mathrm{eff}}_{ijkl}}
\frac{\partial u_k}{\partial x_l}\,\frac{\partial v_i}{\partial x_j}\,d\Omega,
\end{equation}
from which the effective stiffness in \cref{eq:cosserat_ceff} is read off.
In code (\texttt{materials.py}):
\begin{lstlisting}[language=Python]
Cnew = jnp.einsum("jJ,lL,iJkL->ijkl", F, F, C0) / J
\end{lstlisting}
The resulting $\bm C^{\mathrm{eff}}$ generally preserves \emph{major} symmetries but breaks \emph{minor} symmetries, so that
$C^{\mathrm{eff}}_{ijkl} \neq C^{\mathrm{eff}}_{jikl}$ and/or $C^{\mathrm{eff}}_{ijkl} \neq C^{\mathrm{eff}}_{ijlk}$.
This is the signature of Cosserat/micropolar-type behavior when expressed as a single fourth-order tensor.

\section{Triangular carpet transformation (project geometry)}
For the project triangular carpet cloak (piecewise linear boundaries),
a typical mapping keeps $x_1=X_1$ and maps $X_2\in[0,z_2(X_1)]$ to $x_2\in[z_1(x_1),z_2(x_1)]$ via
\begin{equation}
x_1 = X_1,\qquad
x_2 = \left( \frac{z_2(X_1)-z_1(X_1)}{z_2(X_1)} \right) X_2 + z_1(X_1),
\end{equation}
which yields a lower-triangular $\bm F$ (constant $F_{22}$ and piecewise-constant $F_{21}$ for linear $z_1,z_2$).
This is the transformation used in the current project code to build $\bm F(\bm x)$ and then $\bm C^{\mathrm{eff}}(\bm x)$.

\section{Symmetrization for near-cloaking}
If one wants a standard Cauchy-type symmetric tensor (e.g., to avoid micropolar constitutive laws),
a practical approach is to \emph{symmetrize} the Cosserat stiffness. For the 2D in-plane case expressed in an
``augmented Voigt'' notation (including both $12$ and $21$ shear components), one can enforce the minor symmetries by arithmetic averaging,
chosen such that the elastic energy density for symmetric strains is preserved (near-cloaking rationale).
This is the strategy implemented in the project helper \texttt{symmetrize\_stiffness}.

\chapter{Cosserat elasticity (non-symmetric stress)}
\section{Why Cosserat appears in transformation elastodynamics}
In the Cosserat gauge (\cref{eq:cosserat_ceff}), the transformed PDE is written using $\nabla \bm u$ rather than $\bm\varepsilon(\bm u)$.
As a result, the stress-like object $\bm\sigma = \bm C^{\mathrm{eff}} : \nabla \bm u$ is not required to be symmetric.
This aligns with Cosserat (micropolar) continua, where force stresses can be non-symmetric and couple-stresses may appear in more general formulations.

\section{Practical consequences for FEM}
For a standard displacement-based FEM discretization, the main practical differences are:
\begin{itemize}
\item The constitutive mapping is applied to the full gradient $\nabla \bm u$, not only to the symmetric strain.
\item The stiffness tensor must be stored/used in a representation that can encode the broken minor symmetries.
\item If one replaces $\bm C^{\mathrm{eff}}$ with a symmetrized tensor, the model becomes standard Cauchy elasticity (approximate cloak).
\end{itemize}

\chapter{Inverse design options}
This project supports two broad inverse-design routes:
\begin{enumerate}
\item \textbf{Gradient-free optimization} (robust to non-smoothness, expensive per evaluation).
\item \textbf{Gradient-based optimization} (efficient when end-to-end differentiation is stable).
\end{enumerate}

\section{Gradient-free optimization}
\subsection{Bayesian optimization}
Bayesian optimization (BO) models the objective $L(\bm\theta)$ using a probabilistic surrogate (often a Gaussian process),
and selects the next evaluation using an acquisition function (EI, UCB, etc.).
This is attractive when each forward solve is expensive and the parameter dimension is moderate.

\subsection{Bound-constrained L-BFGS-B (quasi-Newton)}
When gradients are available or can be approximated and parameters have bounds,
L-BFGS-B is a strong baseline. It is a limited-memory quasi-Newton method using an active-set strategy for box constraints.
In practice (e.g., \texttt{scipy.optimize.minimize(method='L-BFGS-B')}) it is often competitive even with approximate gradients.

\section{Gradient-based optimization}
\subsection{End-to-end differentiation}
Because \texttt{JAX-FEM} is compatible with automatic differentiation (AD),
one can differentiate a loss $L$ with respect to:
\begin{itemize}
\item geometric parameters of the transformation,
\item constitutive parameters (e.g., tensor field parameters, NN weights),
\item source parameters, etc.
\end{itemize}
The key requirement is that the solver path is differentiable (linear solve + assembly + boundary operators).
For frequency-domain problems, one typically differentiates through the solution map $\bm u(\bm\theta)$ using
either implicit differentiation or AD through the linear solver, depending on the implementation.

\subsection{Stability notes}
For cloaking, the loss is often sensitive to:
(i) boundary reflections (hence good absorbing layers),
(ii) conditioning of the stiffness/mass matrices (anisotropy and strong contrasts),
(iii) discretization error near singular transformations (mitigate with offsets / regularization).

\chapter{Weak formulation for Rayleigh waves in JAX-FEM}
\section{Governing PDE (frequency domain)}
We consider 2D in-plane elastodynamics (plane strain) in the physical domain $\Omega$ with coordinates $\bm x=(x_1,x_2)$.
A typical frequency-domain form (time-harmonic $e^{-i\omega t}$) is:
\begin{equation}
\nabla \cdot \bm\sigma(\bm u) + \omega^2 \rho\, \bm u = \bm f,
\end{equation}
where $\bm\sigma$ is built either from:
\begin{itemize}
\item Cauchy elasticity: $\bm\sigma = \bm C : \bm\varepsilon(\bm u)$, or
\item Cosserat gauge: $\bm\sigma = \bm C^{\mathrm{eff}} : \nabla \bm u$.
\end{itemize}

\section{Weak form}
Let $\bm v$ be an admissible test function (respecting Dirichlet constraints).
The weak form is:
\begin{equation}
\int_{\Omega} (\nabla \bm v) : \bm C^{\mathrm{eff}} : (\nabla \bm u)\, d\Omega
- \omega^2 \int_{\Omega} \rho^{\mathrm{eff}}\, \bm v \cdot \bm u\, d\Omega
= \int_{\Omega} \bm v\cdot \bm f\, d\Omega + \int_{\Gamma_N} \bm v \cdot \bm t\, d\Gamma,
\label{eq:weak}
\end{equation}
where $\bm t$ is prescribed traction on a Neumann boundary $\Gamma_N$.
For Rayleigh waves, the free surface typically satisfies traction-free BC:
$\bm t=\bm 0$.

\section{Mapping to JAX-FEM hooks}
In \texttt{JAX-FEM}, a problem class typically provides maps for:
\begin{itemize}
\item \textbf{Tensor map}: returns the constitutive tensor field (here $\bm C^{\mathrm{eff}}(\bm x)$ or symmetrized).
\item \textbf{Mass map}: returns density-like terms (here $\rho^{\mathrm{eff}}(\bm x)$).
\item \textbf{Surface maps}: implement Neumann/Robin-type integrals, including absorbing boundary contributions.
\end{itemize}
In the project, these are implemented as:
\begin{lstlisting}[language=Python]
def get_tensor_map(): ...
def get_mass_map(): ...
def get_surface_maps(): ...
\end{lstlisting}
and used by the \texttt{RayleighCloakProblem} class.

\chapter{Absorbing boundary conditions (based on absorbing layers)}
\section{Why not PML here}
Perfectly matched layers (PML) for elastodynamics often rely on complex coordinate stretching,
leading to complex-valued systems. If the framework/toolchain assumes real arithmetic, an alternative is needed.

\section{Rayleigh/Caughey damping absorbing layers}
A simple real-valued absorbing layer strategy is to add Rayleigh (Caughey) damping in a buffer region near truncation boundaries:
\begin{equation}
\bm C'(\bm x) = (1+i\eta(\bm x)) \bm C(\bm x), \qquad
\bm \rho'(\bm x) = (1+i\eta_\rho(\bm x)) \rho(\bm x),
\end{equation}
or, equivalently in time domain, introduce damping matrix $\bm C_d = \alpha \bm M + \beta \bm K$.
In frequency domain, Rayleigh damping leads to a modified operator:
\begin{equation}
\left( \bm K + i\omega(\alpha \bm M + \beta \bm K) - \omega^2 \bm M \right)\bm u = \bm f.
\end{equation}
The project adopts a \emph{layer profile} $\xi(\bm x)$ that increases from 0 (interior) to 1 (outer boundary),
and applies damping only in the absorbing layer. This is described and implemented following Semblat et al.

\section{Implementation sketch (layer profile and boundary hooks)}
A typical implementation uses:
\begin{itemize}
\item a scalar profile $\xi(\bm x)\in[0,1]$ in the absorbing region,
\item effective $\alpha(\bm x)=\alpha_{\max}\xi(\bm x)$, $\beta(\bm x)=\beta_{\max}\xi(\bm x)$,
\item assembly contributions to the residual/Jacobian consistent with \cref{eq:weak}.
\end{itemize}
In the project code, the key helper is \texttt{\_xi\_profile} and the absorbing contributions are wired through surface maps.

\chapter{Reference field computation}
\section{Purpose}
A \emph{reference} (target) field is needed to define cloaking losses. Common choices:
\begin{itemize}
\item \textbf{Pristine medium} (no defect, no cloak): reference Rayleigh wave propagation.
\item \textbf{Ideal cloak} (Cosserat, non-symmetric tensor): a ``best possible'' reference if available.
\end{itemize}
In the current project, the baseline is typically the pristine medium response, compared against the cloaked configuration.

\begin{figure}[h]
  \centering
  \includegraphics[width=0.85\textwidth]{figures/reference_field.png}
  \caption{Reference field snapshot (placeholder). Replace with an exported displacement magnitude or component field.}
  \label{fig:ref-field}
\end{figure}

\section{Practical notes}
For stable loss evaluation:
\begin{itemize}
\item record reference on the same mesh and at the same frequency $\omega$,
\item store displacement components and derived quantities (magnitude, surface trace),
\item use consistent normalization (e.g., normalize by reference energy on observation set).
\end{itemize}

\chapter{Loss construction for cloaking}
\section{Observation sets}
Define an observation set $\Gamma_{\mathrm{obs}}$ (often a line/strip on the free surface \emph{after} the cloak)
and optionally a bulk observation region $\Omega_{\mathrm{obs}}$ excluding the cloak/defect.

\section{Canonical loss terms}
Let $\bm u(\bm x)$ be the field for a candidate cloak, and $\bm u_{\mathrm{ref}}(\bm x)$ the reference.
Typical losses include:

\paragraph{(1) Field matching on an observation boundary}
\begin{equation}
L_{\Gamma} = \frac{\int_{\Gamma_{\mathrm{obs}}} \|\bm u - \bm u_{\mathrm{ref}}\|^2\, d\Gamma}{\int_{\Gamma_{\mathrm{obs}}} \|\bm u_{\mathrm{ref}}\|^2\, d\Gamma + \varepsilon}.
\end{equation}

\paragraph{(2) Bulk matching away from the cloak}
\begin{equation}
L_{\Omega} = \frac{\int_{\Omega_{\mathrm{obs}}} \|\bm u - \bm u_{\mathrm{ref}}\|^2\, d\Omega}{\int_{\Omega_{\mathrm{obs}}} \|\bm u_{\mathrm{ref}}\|^2\, d\Omega + \varepsilon}.
\end{equation}

\paragraph{(3) Penalties / regularization}
Regularize parameters $\bm\theta$ (e.g., smoothness of tensor fields, bounds on anisotropy, etc.):
\begin{equation}
L_{\mathrm{reg}} = \lambda_s \|\nabla \bm\theta\|^2 + \lambda_b\,\mathrm{barrier}(\bm\theta).
\end{equation}

\paragraph{Full loss}
\begin{equation}
L(\bm\theta) = w_\Gamma L_{\Gamma} + w_\Omega L_{\Omega} + w_{\mathrm{reg}} L_{\mathrm{reg}}.
\end{equation}

\section{Connection to project implementation}
The current project defines a callable \texttt{cloaking\_loss(...)} that computes a scalar loss based on:
(i) a chosen observation set on the free surface,
(ii) reference field data,
(iii) optional weighting/normalization.

\chapter{Sources (selected)}
\vspace{-0.5em}
\begin{thebibliography}{99}

\bibitem{Milton2006}
G.~W. Milton, M.~Briane, and J.~R. Willis,
``On cloaking for elasticity and physical equations with a transformation invariant form,''
\emph{New Journal of Physics}, vol.~8, p.~248, 2006. doi:10.1088/1367-2630/8/10/248.

\bibitem{NorrisShuvalov2011}
A.~N. Norris and A.~L. Shuvalov,
``Elastic cloaking theory,''
\emph{Wave Motion}, vol.~48, no.~6, pp.~525--538, 2011. doi:10.1016/j.wavemoti.2011.03.002.

\bibitem{Chatzopoulos2023}
Z.~Chatzopoulos, A.~Palermo, A.~Diatta, S.~Guenneau, and A.~Marzani,
``Cloaking Rayleigh waves via symmetrized elastic tensors,''
\emph{International Journal of Engineering Science}, vol.~191, 103899, 2023. doi:10.1016/j.ijengsci.2023.103899.

\bibitem{Semblat2010}
J.-F. Semblat, A.~Gandomzadeh, and L.~Lenti,
``A simple numerical absorbing layer method in elastodynamics,''
\emph{Comptes Rendus M\'ecanique}, vol.~338, no.~1, pp.~24--32, 2010. doi:10.1016/j.crme.2009.12.004.

\bibitem{Shahriari2016}
B.~Shahriari, K.~Swersky, Z.~Wang, R.~P. Adams, and N.~de~Freitas,
``Taking the human out of the loop: A review of Bayesian optimization,''
\emph{Proceedings of the IEEE}, vol.~104, no.~1, pp.~148--175, 2016. doi:10.1109/JPROC.2015.2494218.

\bibitem{Byrd1995}
R.~H. Byrd, P.~Lu, J.~Nocedal, and C.~Zhu,
``A limited memory algorithm for bound constrained optimization,''
\emph{SIAM Journal on Scientific Computing}, vol.~16, no.~5, pp.~1190--1208, 1995. doi:10.1137/0916069.

\bibitem{JAXFEM}
\emph{JAX-FEM: Differentiable Finite Element Method in JAX}, GitHub repository (deepmodeling/jax-fem).\\
\url{https://github.com/deepmodeling/jax-fem}

\bibitem{SciPyLBFGSB}
\emph{SciPy documentation: L-BFGS-B in \texttt{scipy.optimize.minimize}}, SciPy user guide.\\
\url{https://docs.scipy.org/doc/scipy/reference/optimize.minimize-lbfgsb.html}

\end{thebibliography}

\end{document}
